\documentclass[11pt]{article}
\usepackage{amsmath,amssymb}
\usepackage{amsthm}
\usepackage{mathtools}
\usepackage[margin=1in]{geometry}
\usepackage{booktabs}
\usepackage{hyperref}
\usepackage{xcolor}
\hypersetup{colorlinks=true, linkcolor=blue!50!black, urlcolor=blue!50!black}

\newcommand{\sat}[2]{[#1]\!:\!#2^{\infty}}
\newcommand{\satalg}[2]{(#1)\!:\!#2^{\infty}}
\newcommand{\code}[1]{\texttt{#1}}
\newcommand{\ld}{\operatorname{ld}}
\newcommand{\sep}{\operatorname{sep}}
\newcommand{\ini}{\operatorname{ini}}
\newcommand{\si}{\operatorname{si}}
\newcommand{\NF}{\operatorname{NF}}
\newcommand{\dZero}{\mathrm{d\text{-}Zero}}
\newcommand{\dprem}{\mathrm{d\text{-}prem}}
\newcommand{\R}{\mathfrak{R}}
\newcommand{\PP}{\mathbb{P}}
\newcommand{\TT}{\mathbb{T}}
\newcommand{\UU}{\mathbb{U}}
\newcommand{\imp}{\;\Longrightarrow\;}

\title{Coherent, regular, simple, passive:\\
       a reconciliation of the two decomposition traditions\\
       {\large (Rosenfeld--Gr\"obner vs.\ differential Thomas)}}
\author{}
\date{\today}

\begin{document}
\maketitle

\begin{abstract}
\noindent
A companion note to \emph{A New Method of Solving PDEs}, recording the exact
relationship between four notions that recur whenever one compares the
Rosenfeld--Gr\"obner (RG) and differential-Thomas decompositions: a differential
triangular system may be \emph{coherent}, \emph{regular}, \emph{(differentially)
simple}, or \emph{passive} (\emph{= Janet-involutive}). These are \emph{not}
synonyms and they are \emph{not} mutually equivalent. They form a strict
one-directional hierarchy
\[
  \text{simple}\imp\text{regular}\imp\text{coherent},
  \qquad\qquad
  \text{passive}\imp\text{coherent},
\]
with coherence as the common base property (``all $\Delta$-polynomials reduce to
zero'') that makes Rosenfeld's Lemma fire. The note pins each node to a primary
source, quotes the defining conditions verbatim, and flags one correction to an
earlier informal claim (that Seiler's monograph states a coherence
$\Leftrightarrow$ involution \emph{equivalence}): the sourced statement is the
\emph{implication} passive $\Rightarrow$ coherent (Li--Wang 1999, Thm.~3;
Gerdt 1999, Thm.~16; B\"achler et al.\ 2012). It also proves, in
\S\ref{sec:simple-coherent}, that a differentially simple system in
B\"achler et al.'s sense (whose definition, unlike Li--Wang's, does
\emph{not} assume coherence) is coherent --- the arrow a Thomas-based argument
actually needs, and one that avoids importing Wu's notion of passivity.
\end{abstract}

\section{Why this note exists}

RG and differential Thomas both decompose the radical differential ideal
$\{\PP\}$ generated by a finite set $\PP$ of differential polynomials, but they
return components of different \emph{kinds}:
\begin{itemize}
  \item RG returns \emph{regular} differential systems and represents
        $\{\PP\}$ as an \emph{intersection} of regular differential ideals
        (components may overlap);
  \item differential Thomas returns \emph{simple} differential systems and
        gives a \emph{disjoint} decomposition (the solution sets partition).
\end{itemize}

\paragraph{A duality caveat.} The two bullets above are stated at
\emph{different} levels, and it is worth being explicit about it, because
``intersection'' and ``union'' are dual here. RG's ``intersection'' is at the
\emph{ideal} level, $\sqrt{[\PP]}=\bigcap_k\sqrt{[A_k]\!:\!H_k^\infty}$; dually,
at the \emph{variety} level it is a \emph{union}
$\dZero(\PP)=\bigcup_k\dZero(\text{regular}_k)$ of the component solution sets.
Thomas's ``disjoint decomposition'' is stated at the \emph{variety} level. So
the honest common frame is: \emph{both} are unions of solution sets; RG's is a
possibly-\emph{overlapping} union (equivalently, an ideal intersection), while
Thomas's is a \emph{disjoint} union (a partition). The ``intersection vs.\
disjoint-union'' contrast is really ``ideal-intersection (RG) vs.\
variety-partition (Thomas)'', which compares across the duality; each statement
is individually the conventional one for its algorithm.

\paragraph{Where parametric RG sits: a union of intersections.} A third variant,
\emph{parametric} RG (Fakouri--Rahmany--Basiri 2018; the algorithm ported in
\code{\char`~/parametric\_rg}), adds a \emph{disjoint} stratification of parameter
space on top of an RG decomposition of the differential fibre. Writing the
parameter space as a partition into cells $\Omega_{\text{par}}=\bigsqcup_j C_j$,
\[
  \dZero(\PP)\;=\;\bigsqcup_{j}\ \Bigl(\ \underbrace{\textstyle\bigcup_k
      \dZero(\text{regular}_{j,k})}_{\text{RG decomposition over cell }C_j}\ \Bigr),
\]
or, naming each cell's fibre by its ideal representation, over $C_j$ it is
$\bigcap_k\sqrt{[A_{j,k}]\!:\!H^\infty}$. In words:
\[
  \boxed{\ \text{parametric RG}\ =\ \text{a \emph{disjoint} union (over parameter
      cells) of RG intersections (per cell)}\ }
\]
--- the outer union disjoint (Thomas-like, \emph{on the parameter axis}), the
inner object an ordinary RG intersection of regular components (RG-like,
\emph{on the differential fibre}, and possibly overlapping as varieties). It is
a union \emph{of} intersections, never an intersection of unions: across
disjoint parameter cells one can only case-split, never intersect. This is the
precise sense of the ``hybrid'' character noted in \S7 --- disjoint on
parameters, overlapping-regular on the fibre.

Underneath all of these sits a third notion, \emph{coherence} (the condition in
Rosenfeld's Lemma), and on the Janet--Riquier side a fourth, \emph{passivity /
involution}. In loose usage these four words get treated as interchangeable.
They are not. The purpose of this note is to fix the definitions from primary
sources and record the exact implication lattice, so the distinction need not be
re-derived. The single most on-point reference is Li \& Wang (1999), whose very
title is the question --- ``\emph{Coherent, regular and simple systems in zero
decompositions of partial differential systems}'' --- and whose abstract states
its purpose as: \emph{``Rosenfeld's lemma is extended to a more general setting;
relations between passivity and coherence are clarified; regular systems and
simple systems are generalized and proposed.''}

\section{Standing notation}

Fix a differential field $K$ of characteristic $0$ with commuting derivations,
and the differential polynomial ring $\R=K\{X_1,\dots,X_n\}$. For a differential
polynomial $P\in\R\setminus K$: $\ld(P)$ is its \emph{lead} (highest-ranked
derivative under a fixed admissible ranking), $\ini(P)$ its \emph{initial}
(leading coefficient in $\ld(P)$), and $\sep(P)=\partial P/\partial\ld(P)$ its
\emph{separant}. Write $\si(\PP)$ for the set of separants and initials of
$\PP$, and $[\PP]$, $\{\PP\}$ for the differential and radical differential
ideals generated by $\PP$. Saturation: $\sat{\PP}{\,\si(\PP)}$ denotes the
differential ideal quotient by the multiplicative set generated by $\si(\PP)$.
$\dZero(\PP/\mathbb{Q})$ is the differential zero set of $\PP$ off the zero set
of $\mathbb{Q}$; $\dprem$ is the differential pseudo-remainder.

We use ``d-tri set'' for a differentially triangular set (distinct leads,
pairwise partially reduced) and ``d-asc set'' for a differential ascending
(autoreduced) set. Every d-asc set is a d-tri set; the converse fails, and Li \&
Wang work with the weaker d-tri notion throughout.

\section{The four definitions, from the sources}

\subsection{Coherence (Li--Wang 1999, Def.~5)}

The $\Delta$-polynomial of $F,G\in\R\setminus K$ of the same class is the
differential analogue of an $S$-polynomial (construction due to Rosenfeld):
\begin{equation}\label{eq:delta}
  \Delta(F,G)=\sep(G)\,\phi_1 F-\sep(F)\,\phi_2 G,
  \qquad \ld(\phi_1 F)=\ld(\phi_2 G),
\end{equation}
where $\phi_1,\phi_2$ are proper derivative operators of \emph{lowest} order
bringing the two leads to a common derivative.

\begin{quote}
\textbf{Definition (coherent).} A weakly-reduced set $\PP$ is \emph{coherent}
if, for any $F,G\in\PP$ of the same class and any $\Delta(F,G)$ as in
\eqref{eq:delta},
\[
  J\,\Delta(F,G)=C_1\theta_1 P_1+\cdots+C_r\theta_r P_r,
  \qquad \theta_i P_i<\ld(\phi_1 F),
\]
for some SI-product $J$, coefficients $C_i\in\R$, and $P_i\in\PP$.
\end{quote}

In words: \emph{every $\Delta$-polynomial reduces to a combination of
lower-lead derivatives of members of $\PP$}. This is precisely the differential
Buchberger criterion. (``Weakly reduced'' (Def.~3) means every $\ini(P)$ and
$\sep(P)$ is partially reduced w.r.t.\ $\PP$; every d-tri set is weakly reduced.)

\paragraph{What coherence buys: the generalized Rosenfeld Lemma.}
\emph{Li--Wang Theorem 2.} If $\PP\subset\R\setminus K$ is weakly reduced and
coherent, and $F\in\sat{\PP}{\,\si(\PP)}$ is partially reduced w.r.t.\ $\PP$,
then $F\in\satalg{\PP}{\,\si(\PP)}$. That is: coherence is exactly the
hypothesis under which \emph{differential} ideal membership collapses to
\emph{algebraic} ideal membership. This is the theoretical engine of the entire
RG algorithm (Boulier--Lazard--Ollivier--Petitot).

\subsection{Regular (Li--Wang 1999, Def.~6)}

First, the system-level upgrade of coherence: a d-tri system $\langle
\TT,\UU\rangle$ is \emph{coherent} if $\TT$ is coherent and each d-pol in $\UU$
is partially reduced w.r.t.\ $\TT$. Then:

\begin{quote}
\textbf{Definition (regular).} A d-tri system $\langle\TT,\UU\rangle$ is
\emph{regular} if it is a coherent \emph{and} elementarily-triangular (e-tri)
system.
\[
  \boxed{\ \text{regular}\ =\ \text{coherent}\ +\ \text{e-tri}\ }
\]
\end{quote}

Here $\langle\TT,\UU\rangle$ is \emph{e-tri} (Def.~1) if \emph{every} point in
$\mathrm{Zero}(\TT/\UU)$ --- the \emph{algebraic}, not just differential, zero
set --- fails to annihilate any element of $\si(\TT)$; equivalently the initials
and separants are invertible on the whole (algebraic) component, not merely
generically. Li \& Wang note this \emph{generalizes} Boulier et al.'s regular
chains in two ways: $\TT$ need not be a d-asc set, and $\si(\TT)$ need not be
contained in $\UU$.

\paragraph{What regularity buys.} \emph{Li--Wang Theorem 8.} For a regular
$\langle\TT,\UU\rangle$: a d-pol $F$ lies in $\sat{\TT}{\,\UU}$ iff a partial
pseudo-remainder of $F$ w.r.t.\ $\TT$ lies in $\satalg{\TT}{\,\UU}$; and both
$\sat{\TT}{\,\UU}$ and $\satalg{\TT}{\,\UU}$ are \emph{radical}. Regular systems
are the components of the RG ``weak form''
$\dZero(\PP/\mathbb{Q})=\bigcup_i\dZero(\TT_i/\UU_i)$ --- an
\emph{intersection}-style decomposition whose components can overlap.

\subsection{Differentially simple (Li--Wang 1999, Def.~8; algebraic core Def.~2)}

The algebraic prototype (Def.~2, after Thomas and Wang): a triangular system
$\langle\TT,\UU\rangle$ in $K[\mathbf x]$ is \emph{w-simple} if for each
$1\le k\le n$ it is triangular, has non-vanishing initials on the lower solution
set, and is \emph{square-free} (the univariate image $T(\bar{\mathbf
x}^{\{p-1\}},x_p)$ is square-free). Lifting to the differential setting:

\begin{quote}
\textbf{Definition (differentially simple, d-sim).} A \emph{coherent} d-tri
system $\langle\TT,\UU\rangle$ in $\R$ is \emph{differentially simple} if, when
considered as a triangular system in the underlying (non-differential) ring, it
is w-simple.
\end{quote}

So, informally,
\[
  \text{simple}\ =\ \text{regular}\ +\ \text{square-free / w-simple conditions}.
\]
The disjoint decomposition $\dZero(\PP/\mathbb{Q})=\bigcup_i\dZero(\TT_i/\UU_i)$
into d-sim systems is the (differential) \emph{Thomas decomposition}; its
components \emph{partition} the solution set.

\paragraph{What simplicity buys.} \emph{Li--Wang Theorem 9 \& Corollary 2.} A
d-sim system has $\dZero(\TT/\UU)\ne\emptyset$, and for any d-pol $P$,
\[
  \dZero(\TT/\UU)\subset\dZero(P)\iff\dprem(P,\TT)=0 .
\]
Moreover $\TT$ is then a characteristic set (in Ritt's sense) of the prime
differential ideal $\sat{\TT}{\,\UU}$. And, top of Li--Wang p.~58: ``\emph{a
d-sim system must be regular}.''

\subsection{Passive / involutive (Gerdt 1999; used by differential Thomas)}

``Passive'' is the Janet--Riquier word for the same completeness idea, phrased
via an \emph{involutive division} (Janet division for the Thomas construction).
Gerdt states the terminology bridge in prose: \emph{``A system satisfying the
Janet involutivity conditions is often called \textbf{passive}.''} The condition
itself:

\begin{quote}
\textbf{Definition (involutive, Gerdt 1999 Def.~14).} An $L$-autoreduced set $F$
is \emph{$L$-involutive} if
\[
  (\forall f\in F)\,(\forall\alpha\in\mathbb{N}^n)\;
  \bigl[\,\NF_L(\partial_\alpha f,\,F)=0\,\bigr].
\]
\textbf{Local criterion (Gerdt 1999 Thm.~16).} $F$ is $L$-involutive
(for a continuous involutive division $L$) iff
\[
  (\forall f\in F)\,(\forall x_i\in\mathrm{NM}_L(f,F))\;
  \bigl[\,\NF_L(\partial_{x_i}\!\cdot f,\,F)=0\,\bigr],
\]
i.e.\ only the \emph{non-multiplicative} prolongations need reduce to zero.
\end{quote}

This finite, checkable criterion is ``the involutivity criterion'' cited (in its
nonlinear-PDE form, Gerdt 2008) by B\"achler et al.\ (2012). Note the exact
match with the involutive clause in the B\"achler et al.\ definition of a
differential simple system: \emph{``$S$ is (Janet) involutive if all
non-admissible prolongations in $(S_T)^=$ reduce to zero by $(S_T)^=$.''} Gerdt
also records \emph{Theorem 17}: an $L$-involutive basis of $[F]$ is a
\emph{differential Gr\"obner basis}; and \emph{Theorem 18}, an involutive
analogue of the Buchberger chain criterion --- the same $S$-polynomial /
$\Delta$-polynomial content that, on the RG side, \emph{is} coherence.

\section{The connecting theorem: passive $\Rightarrow$ coherent}

The bridge between the two vocabularies is a single, one-directional theorem,
stated and proved in Li--Wang and echoed in the Thomas literature.

\begin{quote}
\textbf{Li--Wang 1999, Theorem 3.} \emph{A passive d-asc set is coherent.}
\end{quote}

\noindent
The proof (Li--Wang \S4.1, ``Passivity and coherence'') runs through Wu's
tuple-decomposition and passivity theorems; the section's framing is explicit
that this direction is the useful one: ``\emph{coherent d-asc sets are, in some
sense, simpler than passive d-asc sets}.'' The converse is \emph{not} claimed.

\paragraph{Whose ``passive''?} Li--Wang do \emph{not} define the term. \S4.1
opens, immediately above Theorem~3:

\begin{quote}
``Passive systems were investigated by Ritt$^{[3]}$. \textbf{The reader is
referred to Wu ([2]) for a general theory for passive d-asc sets and precise
definitions of terms not defined in this section.}''
\end{quote}

\noindent
--- reference [2] being Wu, \emph{On the foundation of algebraic differential
geometry}, Sys.\ Sci.\ Math.\ Scis.\ \textbf{2} (1989) 289--312. The proof of
Theorem~3 is correspondingly conducted in imported Wu machinery
($\mathrm{CTUP}_k(\mathbb{A})$, the tuple-decomposition theorem, the passivity
theorem, M-derivatives, non-multiplier coordinates), none of it defined in
Li--Wang.

So Theorem~3 is a statement about \emph{Wu-passive} d-asc sets. It is
\emph{not}, as stated, a statement about Janet-involutivity in the sense of
Gerdt Def.~14 or of B\"achler et al.\ Def.~3.5. The two notions are surely
closely related --- both formalize Janet--Riquier completion --- but no source
consulted here states the identification, and this note should not assume it.

\noindent
For the Janet-involutive notion the implication has its own, same-framework
source: B\"achler, Gerdt, Lange-Hegermann \& Robertz (2012), \S1: ``\emph{the
output systems are Janet-involutive $\dots$ and \textbf{hence they are
coherent}},'' and, on the algorithm, ``\emph{the involutive approach, which we
utilize to make the subsystems coherent}.'' Direction: involutive
$\Rightarrow$ coherent, via Gerdt's involutivity criterion.

\noindent
\textbf{Practical upshot.} There are two arrows into coherence, with different
antecedents and different sources. Work that starts from a differential Thomas
decomposition should cite B\"achler et al.\ \S1, not Li--Wang Theorem~3; the
latter requires Wu's definition of passivity and an identification nobody has
written down.

\paragraph{Correction to an earlier informal claim.} It is tempting (and was
asserted in passing during the discussion that produced this note) to attribute
a coherence $\Leftrightarrow$ involution \emph{equivalence} to W.~Seiler's
monograph \emph{Involution} (2010) and, worse, to his two open-access AAECC 2009
papers (``A combinatorial approach to involution and $\delta$-regularity, I \&
II''). This is not supported by those two papers: they are commutative-algebra
papers on involutive/Pommaret bases and $\delta$-regularity and contain
\emph{zero} occurrences of ``coherence,'' ``Rosenfeld,'' or ``regular
differential.'' They connect involution to Janet--Riquier \emph{passivity}, not
to RG coherence. The defensible, sourced statement is the \emph{implication}
above, not an equivalence; for Janet-involutivity its home is B\"achler et al.\
(2012), \S1, resting on Gerdt 1999/2008, and for Wu-passivity it is Li--Wang
1999 (Thm.~3), resting on Wu 1989 --- in neither case Seiler. (Seiler's $\delta$-regularity work is nonetheless
the right reference for \emph{why} Pommaret-involution --- unlike Janet-involution
and unlike coherence --- is coordinate-dependent, which is a second, independent
reason the two notions cannot be identical.)

\section{A second route to coherence: Thomas-simple $\Rightarrow$ coherent}
\label{sec:simple-coherent}

\S4 leaves the Janet-involutive arrow into coherence resting on an assertion
(B\"achler et al.\ \S1) rather than a proof. This section supplies a proof of a
\emph{weaker but sufficient} statement --- one whose antecedent is the full
definition of a differentially simple system rather than involutivity alone ---
which is what any argument starting from a differential Thomas decomposition
actually has in hand.

\paragraph{A terminological trap first.} Two different conditions are called
``simple'' in this literature, and only one of them makes the implication
non-trivial.

\begin{itemize}
  \item \textbf{Li--Wang d-sim} (Def.~8; \S3.3 above) is defined as a
        \emph{coherent} d-tri system that is w-simple. Coherence sits
        \emph{inside} the definition, so ``simple $\Rightarrow$ coherent'' is a
        tautology for this notion and the arrow in the lattice below records
        nothing mathematical.
  \item \textbf{BGLR simple} (Def.~3.3) is (S1) algebraically simple ---
        triangular, non-vanishing initials, square-free --- plus (S2)
        Janet-involutive, (S3) $S^=$ minimal, (S4) no inequation reducible
        modulo $S^=$. \textbf{Coherence is not among these conditions.} For this
        notion ``simple $\Rightarrow$ coherent'' is a theorem.
\end{itemize}

\noindent
The differential Thomas algorithm outputs BGLR-simple systems, so it is the
second notion that is available downstream, and the second implication that has
to be established.

\paragraph{Coherence in BLOP's formulation.} Li--Wang Def.~5 (\S3.1 above) is
stated for a set $\PP$; the equivalent system-level phrasing in
Boulier--Lazard--Ollivier--Petitot is the one that meshes with the Thomas-side
results, so we use it here. For a set $A$ of differential polynomials and a
derivative $v$, write $A_v=\{\theta p \mid p\in A,\ \theta\in\Theta,\
\ld\theta p\le v\}$ and $\R_v$ for the subring on derivatives of rank $\le v$.

\begin{itemize}
  \item \textbf{(Def.~17, critical pair.)} $\{p_1,p_2\}$ is a \emph{critical
        pair} if $\ld p_1$ and $\ld p_2$ have a common derivative. If $\ld p_2$
        is a derivative of $\ld p_1$ the pair is a \emph{reduction} critical
        pair.
  \item \textbf{(Def.~18, $\Delta$-polynomial.)} With
        $\theta_iu=\ld p_i$, $\theta_{12}u=\mathrm{lcd}(\theta_1u,\theta_2u)$ and
        $\operatorname{rank}p_1<\operatorname{rank}p_2$,
        \[
          \Delta(p_1,p_2)=\sep(p_1)\frac{\theta_{12}}{\theta_2}p_2
                         -\sep(p_2)\frac{\theta_{12}}{\theta_1}p_1
        \]
        in the non-reduction case, and
        $\Delta(p_1,p_2)=p_2\ \mathrm{full\text{-}rem}\ (\theta_2/\theta_1)p_1$
        in the reduction case.
  \item \textbf{(Def.~19, solved.)} The pair is \emph{solved by} $A=0$,
        $S\ne 0$ if there is a derivative $v<\theta_{12}u$ with
        $\Delta(p_1,p_2)\in\satalg{A_v}{(S\cap\R_v)}$.
  \item \textbf{(Def.~22, condition C3.)} $A=0$, $S\ne0$ is \emph{coherent} if
        every critical pair of $A$ is solved by it.
\end{itemize}

\paragraph{The imported ingredient.} One result does the work, and it is a
Thomas-side result with no coherence in its hypothesis:

\begin{quote}
\textbf{Robertz, \emph{Formal Algorithmic Elimination for PDEs}, Prop.~2.2.50}
(restated as Prop.~3.31 of the CIRM notes). \emph{Let $S$ be a simple
differential system with equations $p_1=0,\dots,p_s=0$, let $E$ be the
differential ideal generated by $p_1,\dots,p_s$, and let $q$ be the product of
the initials and separants of $p_1,\dots,p_s$. Then $E:q^\infty$ is a radical
differential ideal, and for $p\in\R$ we have $p\in E:q^\infty$ if and only if
the pseudo-remainder of $p$ modulo $p_1,\dots,p_s$ and their derivatives is
zero.}
\end{quote}

\begin{quote}
\textbf{BLOP Lemma 20.} \emph{Let $\{p_1,p_2\}$ be a critical pair with
$\ld p_1\ne\ld p_2$, and let $A=0$, $S\ne0$ be a differential system with
$H_A\subset S$. If $\Delta(p_1,p_2)\ \mathrm{full\text{-}rem}\ A=0$ then
$\{p_1,p_2\}$ is solved by $A=0$, $S\ne0$.}
\end{quote}

\paragraph{The statement.}

\begin{quote}
\textbf{Proposition.} Let $S=(S^=,S^{\ne})$ be a differentially simple system in
the sense of BGLR Def.~3.3, with $A=S^==\{p_1,\dots,p_s\}$, and let
$H=H_A=\si(A)$ be its initials and separants. Then every critical pair of $A$ is
solved by $A=0$, $H\ne0$; that is, $A$ is coherent.
\end{quote}

\begin{proof}
Let $\{p_1,p_2\}$ be a critical pair of $A$, with notation as in Def.~18, and
let $E=[A]$, $q=\prod_{p\in A}\ini(p)\sep(p)$.

\emph{Step 1: $\Delta(p_1,p_2)\in E$.} In the non-reduction case each
$(\theta_{12}/\theta_i)p_i$ is a derivative of an element of $A$ and so lies in
$E$; since $E$ is an ideal, the combination
$\sep(p_1)(\theta_{12}/\theta_2)p_2-\sep(p_2)(\theta_{12}/\theta_1)p_1$ lies in
$E$ as well. In the reduction case Ritt full reduction supplies an SI-product
$h$ with $h\,p_2-\Delta(p_1,p_2)\in[p_1]$, so
$\Delta(p_1,p_2)=h\,p_2-(\text{element of }[p_1])\in E$. Either way
$\Delta(p_1,p_2)\in E\subseteq E:q^\infty$.

\emph{Step 2: $\Delta(p_1,p_2)$ pseudo-reduces to zero modulo $A$.} $S$ is
simple, so Robertz Prop.~2.2.50 applies; take its ``only if'' direction with
$p=\Delta(p_1,p_2)$, which lies in $E:q^\infty$ by Step~1.

\emph{Step 3: the pair is solved.} By (S1)/(A1) the elements of $A$ have
pairwise distinct leaders, so $\ld p_1\ne\ld p_2$, and $H_A\subset H$ by
construction; BLOP Lemma~20 applies. Explicitly, put $v=\ld\Delta(p_1,p_2)$,
which is $<\theta_{12}u$, and $\bar A=\{p\in A\mid \operatorname{rank}p\le
\operatorname{rank}\Delta(p_1,p_2)\}$. By the specification of Ritt reduction
there are $h_1,\dots,h_n\in H_{\bar A}$ and exponents $\alpha_i$ with
$h_1^{\alpha_1}\cdots h_n^{\alpha_n}\Delta(p_1,p_2)\in(\bar A_v)$. Since
$H_{\bar A}\subseteq H\cap\R_v$ and $\bar A_v\subseteq A_v$ we get
$\Delta(p_1,p_2)\in\satalg{A_v}{(H\cap\R_v)}$ with $v<\theta_{12}u$, which is
Def.~19.
\end{proof}

\paragraph{On the non-canonicity of Ritt reduction.} Ritt full reduction is not
canonical --- the remainder depends on which reducer is chosen at each step ---
so ``the pseudo-remainder is zero'' needs justification as a well-posed
condition. Robertz's proposition supplies it: it is an \emph{equivalence} with
an ideal membership, so for a simple system the vanishing of the remainder is
independent of the reduction strategy. Step~2 is therefore unambiguous.

\paragraph{No circularity.} The obvious worry is that Robertz proves Prop.~2.2.50
by way of Rosenfeld's Lemma, which takes coherence as a hypothesis --- in which
case the argument above would be circular. It does not: ``coherent'' occurs
nowhere in the body of the CIRM notes, and Rosenfeld appears only in the
introduction (as a contrast to Thomas) and in the bibliography. The route to
2.2.50 is the Thomas/Janet theory together with the Ritt--Raudenbush
Nullstellensatz.

\paragraph{What this settles, and what it does not.} It settles the practical
question raised at the end of \S4: work that starts from a differential Thomas
decomposition does \emph{not} need Li--Wang Thm.~3, and does not need to import
Wu's definition of passivity or assume an unwritten identification of
Wu-passivity with Janet-involutivity. It can derive coherence of the output
systems from a Thomas-side result. It also explains B\"achler et al.\ \S1: read
in context --- ``\emph{the output systems} are Janet-involutive $\dots$ and
hence they are coherent'' --- their claim is about the simple systems their
algorithm returns, not about involutivity in isolation, and in that reading it
is correct.

It does \emph{not} prove Janet-involutive $\Rightarrow$ coherent. Step~2 uses
Prop.~2.2.50, whose hypothesis is (S1)--(S4) together; involutivity enters only
as (S2). The bare implication, with involutivity as the sole antecedent, remains
unlocated in the literature and unproved here.

\section{The implication lattice}

Collecting Definitions 5, 6, 8 and Theorems 2, 3 of Li--Wang with Gerdt's
Theorem 16:
\[
\begin{gathered}
  \underbrace{\text{Janet-involutive}}_{\substack{\text{Gerdt Def.~14/Thm.~16}\\ \text{BGLR Def.~3.5}}}
  \;\xRightarrow{\ \text{BGLR \S1}\ }\;
  \text{coherent}
  \\[1ex]
  \underbrace{\text{Wu-passive}}_{\text{Wu 1989}}
  \;\xRightarrow{\ \text{L--W Thm.~3}\ }\;
  \text{coherent}
  \\[1.5ex]
  \underbrace{\text{simple}}_{\text{L--W Def.~8}}
  \imp
  \underbrace{\text{regular}}_{\text{L--W Def.~6}}
  \;=\;\text{coherent}+\text{e-tri}
  \imp
  \text{coherent}
  \\[1.5ex]
  \underbrace{\text{simple}}_{\text{BGLR Def.~3.3}}
  \;\xRightarrow{\ \S\ref{sec:simple-coherent}\ }\;
  \text{coherent}.
\end{gathered}
\]

\noindent
The last two rows are not the same statement. Li--Wang's d-sim has coherence
inside its definition, so its arrow is a tautology; BGLR's simple does not, so
its arrow is the proposition of \S\ref{sec:simple-coherent}. The first row
remains an assertion (B\"achler et al.\ \S1) rather than a proved implication.

\begin{center}
\begin{tabular}{@{}p{2.3cm}p{5.6cm}p{3.0cm}p{2.7cm}@{}}
\toprule
Node & Defining condition & Decomposition & Source \\
\midrule
coherent & $\Delta$-polys reduce to $0$ & (property) & L--W Def.~5 \\
Janet-involutive & non-mult.\ prolongations reduce to $0$ & (property) & Gerdt Def.~14, Thm.~16; BGLR Def.~3.5 \\
Wu-passive & (not defined in L--W; see Wu 1989) & (property) & Wu 1989 \\
regular & coherent $+$ e-tri (sep./init.\ invertible on the component) & intersection (RG weak form) & L--W Def.~6 \\
simple (d-sim) & \emph{coherent} d-tri $+$ w-simple (triangular, non-vanishing init., square-free) & disjoint (Thomas) & L--W Def.~8 \\
simple (BGLR) & alg.\ simple $+$ involutive $+$ minimal $+$ inequation-reduced; coherence \emph{not} assumed & disjoint (Thomas) & BGLR Def.~3.3; coherence proved \S\ref{sec:simple-coherent} \\
\bottomrule
\end{tabular}
\end{center}

\noindent
Every arrow is \emph{strict} and \emph{one-directional}:
\begin{itemize}
  \item \textbf{coherent} is the base property --- it makes Rosenfeld's Lemma
        fire (differential $\to$ algebraic), and nothing more is assumed.
  \item \textbf{passive $\Rightarrow$ coherent} but not conversely: passivity is
        a specific, division-dependent (Janet) completion; coherence is the
        weaker, division-free, intrinsic property it implies (L--W Thm.~3).
  \item \textbf{regular $\Rightarrow$ coherent} but not conversely: regular adds
        the e-tri (whole-component invertibility of separants/initials) upgrade,
        which is what makes the saturated ideals radical (L--W Thm.~8).
  \item \textbf{simple $\Rightarrow$ regular} but not conversely: simple adds
        square-freeness / the w-simple conditions, which is what buys
        \emph{disjointness} of the Thomas decomposition and the characteristic-set
        property (L--W Cor.~2).
  \item \textbf{BGLR-simple $\Rightarrow$ coherent} is the one arrow here with a
        proof written out rather than cited (\S\ref{sec:simple-coherent}), and
        it is the one a Thomas-based argument needs, since BGLR's definition ---
        unlike Li--Wang's --- does not assume coherence.
\end{itemize}

\section{Consequences for the RG-vs-Thomas comparison}

\begin{itemize}
  \item \textbf{Coherence is not the discriminator between RG and Thomas.} Both
        traditions require it (RG via Rosenfeld's Lemma; Thomas via the
        involutive clause that \emph{implies} it). The real discriminators are
        the \emph{extra} conditions: e-tri (regular) vs.\ e-tri $+$ square-free
        (simple), i.e.\ \emph{overlapping intersection} vs.\ \emph{disjoint
        partition}.
  \item \textbf{$\Delta$-polynomial $=$ differential $S$-polynomial.} B\"achler
        et al.\ (2012) say it outright (``the analogon of $s$-polynomials in
        differential algebra''), and Gerdt's Theorem 18 is the involutive
        Buchberger chain criterion. So the RG/Thomas relationship mirrors the
        algebraic Gr\"obner-basis / involutive-basis relationship: every
        involutive basis is a Gr\"obner basis but generally larger and more
        structured; likewise passive $\Rightarrow$ coherent, and simple is the
        more refined (disjoint) decomposition.
  \item \textbf{On ``is a differential Thomas decomposition always finer than a
        (parametric) RG decomposition?''} Fibrewise (same ranking, differential
        structure), yes: simple $\Rightarrow$ regular, and simple exposes the
        degenerate branches RG saturation folds away. But \emph{globally},
        including a parameter stratification, the two can become
        \emph{incomparable}, because they split along the discriminant loci that
        \emph{their own} orderings produce, and parametric RG carries a
        distinguished parameter layer (with deferred genericity conditions) that
        Thomas --- having ``no notion of parameter'' --- does not reproduce. The
        one-directional lattice here is what holds \emph{after} fixing a common
        ambient ring and ranking; it is not a claim of global refinement.
\end{itemize}

\section{Primary sources}

\begin{itemize}
  \item Z.~Li \& D.~Wang, \emph{Coherent, regular and simple systems in zero
        decompositions of partial differential systems}, Systems Science and
        Mathematical Sciences \textbf{12} (Suppl.), May 1999, 43--60.
        [Defs.~5 (coherent), 6 (regular), 8 (d-sim); Thm.~2 (generalized
        Rosenfeld), Thm.~3 (passive $\Rightarrow$ coherent), Thm.~8 (radicality),
        Thm.~9/Cor.~2 (simple).] Local archive:
        \nolinkurl{~/project/papers/li-wang-1999-coherent-regular-simple-systems-zero-decompositions-pde.pdf}
        (scanned; page images read directly).
  \item V.~P.~Gerdt, \emph{Completion of Linear Differential Systems to
        Involution}, CASC 1999, 115--137 (arXiv:math/9909114). [Def.~14
        (involutive), Thm.~16 (local involutivity criterion: non-multiplicative
        prolongations reduce to $0$), Thm.~17 (involutive $\Rightarrow$
        differential Gr\"obner), Thm.~18 (involutive Buchberger chain criterion);
        prose: passive $=$ Janet-involutive.] Local archive:
        \nolinkurl{~/project/papers/gerdt-1999-completion-linear-differential-systems-to-involution.pdf}.
  \item V.~P.~Gerdt, \emph{On Decomposition of Algebraic PDE Systems into Simple
        Subsystems}, Acta Appl.\ Math.\ \textbf{101} (2008), 39--51. [The
        involutivity criterion in nonlinear-PDE form, as cited by B\"achler et
        al.]
  \item T.~B\"achler, V.~Gerdt, M.~Lange-Hegermann \& D.~Robertz, \emph{Thomas
        Decomposition of Algebraic and Differential Systems}, J.~Symbolic
        Comput.\ \textbf{47} (2012), 1233--1266 (arXiv:1008.3767). [Def.~2.1
        (algebraic simple), Def.~3.3 (differential simple $=$ algebraically
        simple $+$ involutive $+$ minimal $+$ inequation-reduced); ``Janet-involutive
        $\dots$ hence coherent''; $\Delta$-poly $=$ $s$-poly analogue.] Local
        archives (published + preprint) in \nolinkurl{~/project/papers/}.
  \item F.~Boulier, D.~Lazard, F.~Ollivier \& M.~Petitot, \emph{Computing
        representations for radicals of finitely generated differential ideals},
        AAECC \textbf{20} (2009), 73--121. [Def.~13 (differentially triangular),
        Defs.~17--19 (critical pairs, $\Delta$-polynomials, solved pairs),
        Lemma~20 ($\Delta$ reduces to $0$ $\Rightarrow$ pair solved), Def.~22
        (regular differential system; C3 is coherence), Thm.~23 (Rosenfeld's
        Lemma). The system-level phrasing of coherence used in
        \S\ref{sec:simple-coherent}.] Local archive:
        \nolinkurl{~/project/papers/boulier-lazard-ollivier-petitot-2009-radical-differential-ideal.pdf}.
  \item D.~Robertz, \emph{Formal Algorithmic Elimination for PDEs}, Lecture
        Notes in Mathematics \textbf{2121}, Springer, 2014. [Prop.~2.2.50:
        pseudo-reduction modulo a simple differential system decides membership
        in $E:q^\infty$ --- the ingredient of \S\ref{sec:simple-coherent}.
        Restated as Prop.~3.31 of the CIRM notes, \emph{Formal Methods for
        Systems of Partial Differential Equations} (2018), which is freely
        available and contains no appeal to coherence or Rosenfeld's Lemma in
        its development.] Local archives (book appendix~A + both CIRM versions)
        in \nolinkurl{~/project/papers/}.
  \item W.~M.~Seiler, \emph{A Combinatorial Approach to Involution and
        $\delta$-Regularity, I \& II}, AAECC \textbf{20} (2009). [Involutive /
        Pommaret bases, $\delta$-regularity. \emph{Not} a source for a
        coherence $\Leftrightarrow$ involution equivalence --- see \S4
        correction. Relevant for the coordinate-dependence of Pommaret
        involution.]
\end{itemize}

\bigskip
\noindent\emph{Note prepared as computational/literature backing for
NewMethod; reconciles the terminology used across the RG and differential-Thomas
ports (\code{joca.sage}, \code{joca-thomas-native.sage}, and the parametric-RG
port).}

\end{document}
